\documentclass[11pt,a4paper]{article}
\usepackage{DDTA_DERMATRIAGE_GUIDED_REVIEW_STYLE_R1}
\providecommand{\tightlist}{\setlength{\itemsep}{0pt}\setlength{\parskip}{0pt}}
\begin{document}
\selectlanguage{italian}
\begin{titlepage}
\thispagestyle{empty}
\vspace*{10mm}
{\sffamily\bfseries\color{DDTANavy}\fontsize{15}{18}\selectfont Documentation-Driven Threat Analysis}\par
\vspace{5mm}
{\sffamily\bfseries\color{DDTANavy}\fontsize{27}{31}\selectfont DermaTriage Governed DDTA Documentation}\par
\vspace{2mm}
{\sffamily\color{DDTABlue}\fontsize{17}{21}\selectfont A5 Candidate R1 - Pending A6-A12}\par
\vspace{10mm}
\begin{tcolorbox}[enhanced,arc=3pt,boxrule=0pt,colback=DDTANavy,left=12pt,right=12pt,top=10pt,bottom=10pt]
{\color{white}\sffamily\bfseries Current research artifact}\par
{\color{white!88}\small Repository baseline: \texttt{0e60754d21aa24ea487f3f60803b6b0cce8d2e2b}. Methodology authority remains R4 unless this document explicitly states otherwise.}
\end{tcolorbox}
\vfill
{\sffamily\small\color{DDTAGray} 5 settembre 2026}
\end{titlepage}
\tableofcontents
\clearpage
\section{DermaTriage - Governed DDTA Documentation
Candidate}\label{dermatriage---governed-ddta-documentation-candidate}

\textbf{Status:} A5 CLOSED CANDIDATE / NOT YET PROMOTED\\
\textbf{Coverage:} project framing through FunctionalRequirement
cumulative D3 closure\\
\textbf{Pending before final documentation promotion:} A6-A12\\
\textbf{Methodology authority:}
\texttt{DDTA\_DOCUMENTATION\_BA\_AUTHORING\_GUIDE\_R4}\\
\textbf{Repository baseline:}
\texttt{0e60754d21aa24ea487f3f60803b6b0cce8d2e2b}\\
\textbf{Authoritative source package:}
\texttt{DermaTriage-Docs-20260830T152637Z-1-001.zip}\\
\textbf{Source package SHA-256:}
\texttt{E9ED2C507BEFB95F54A52084687CD1E8798863AE81CF69D09568864D8CBF280E}

\begin{quote}
This file is the \textbf{clean project-document candidate}. Research
commentary and methodological reasoning are intentionally excluded.
Source-supported implementation/configuration evidence is retained only
in clearly marked binding/evidence sections. A6-A12 may still refine
split, specialization, security, terminology, bindings and promotion
readiness.
\end{quote}

\subsection{1. Project problem framing}\label{1-project-problem-framing}

DermaTriage addresses the problem of supporting early triage of
dermatological cases concerning potentially oncological skin lesions,
using the information available for the case to determine a priority of
care and support routing toward appropriate specialist assistance.

\subsubsection{Scope boundary}\label{scope-boundary}

Governed in the current candidate:

\begin{itemize}
\tightlist
\item
  urgency/priority-oriented triage behavior;
\item
  specialist-routing responsibility at macro level;
\item
  management of clinical review results;
\item
  controlled adaptation using accumulated clinical review evidence.
\end{itemize}

Not established as governed project authority:

\begin{itemize}
\tightlist
\item
  authority for definitive clinical diagnosis;
\item
  complete specialist-selection/booking responsibility;
\item
  automatic production-deployment authority after adaptation
  qualification;
\item
  complete semantics of every source-level configuration value.
\end{itemize}

\subsection{2. Canonical macro view}\label{2-canonical-macro-view}

\begin{verbatim}
DERMATRiAGE
|
|-- MR-01  Valutazione di triage del caso dermatologico
|
|-- MR-02  Indirizzamento specialistico
|
|-- MR-03  Gestione della validazione clinica degli esiti
|
`-- MR-04  Adattamento controllato sulla base della revisione clinica
           dependsOn MR-03
\end{verbatim}

Stable IDs are identities, not reading order. The narrative projection
above is for human readability.

\section{3. MR-01 - Valutazione di triage del caso
dermatologico}\label{3-mr-01---valutazione-di-triage-del-caso-dermatologico}

\subsection{Intent}\label{intent}

Determinare, dalle informazioni disponibili sul caso dermatologico, una
valutazione di urgenza utile a stabilire la priorita\textquotesingle{}
di presa in carico specialistica.

\subsection{Context and scope}\label{context-and-scope}

The responsibility concerns urgency/priority evaluation for a
dermatological case. It does not establish authority for definitive
clinical diagnosis and does not own the complete specialist-routing
responsibility.

\subsection{DEC-01 - Continuita\textquotesingle{} del triage in assenza
di
immagine}\label{dec-01---continuita-del-triage-in-assenza-di-immagine}

\textbf{Decision kind:} SELECTABLE\\
\textbf{Status:} PASS

DermaTriage permits urgency evaluation to continue when no image is
available, using available symptom information rather than treating
image availability as an absolute precondition for triage.

\subsubsection{FR-01 - Determinazione dell\textquotesingle urgenza su
base sintomatologica in assenza di
immagine}\label{fr-01---determinazione-dellurgenza-su-base-sintomatologica-in-assenza-di-immagine}

\textbf{Parent:} DEC-01\\
\textbf{Status:} PASS

\textbf{Normative clause}

\begin{quote}
Quando per il caso dermatologico non e\textquotesingle{} disponibile
un\textquotesingle immagine, DermaTriage MUST determinare
l\textquotesingle urgenza sulla base delle informazioni sintomatologiche
disponibili per il caso.
\end{quote}

Source examples such as itching, bleeding, growing, changing and pain
illustrate current symptom evidence and are not treated here as a closed
normative vocabulary.

\subsubsection{Preserved gap}\label{preserved-gap}

\texttt{GAP-DERMA-NOIMAGE-INPUT-01} - the complete required/optional
symptom input contract and missing-value semantics are not specified.

\subsection{DEC-02 - Adozione della P-scale per la
priorita\textquotesingle{}
operativa}\label{dec-02---adozione-della-p-scale-per-la-priorita-operativa}

\textbf{Decision kind:} SELECTABLE\\
\textbf{Status:} PASS

DermaTriage represents operational triage priority using P1-P4, derived
from the urgency evaluation of the case.

\subsubsection{FR-02 - Derivazione della priorita\textquotesingle{}
P-scale dalla valutazione di
urgenza}\label{fr-02---derivazione-della-priorita-p-scale-dalla-valutazione-di-urgenza}

\textbf{Parent:} DEC-02\\
\textbf{Status:} PASS

\textbf{Normative clauses}

\begin{verbatim}
HIGH + confidence > 0.85 -> P1
HIGH                     -> P2
MEDIUM                   -> P3
LOW                      -> P4
\end{verbatim}

The strict comparator \texttt{\textgreater{}} is preserved.

\subsubsection{Preserved gaps}\label{preserved-gaps}

\begin{itemize}
\tightlist
\item
  \texttt{GAP-DERMA-PMAP-INPUT-01} - complete semantics of the
  urgency/confidence inputs, including missing/invalid values, are not
  specified.
\item
  \texttt{GAP-DERMA-NOIMAGE-PMAP-BINDING-01} - the documentation does
  not establish that the urgency produced by FR-01 necessarily provides
  every input required by FR-02.
\item
  \texttt{GAP-DERMA-SLA-01} - source values 24h/48h/72h/7 days are
  preserved as evidence, but trigger event, owned outcome,
  responsibility and whether the values are
  limits/targets/recommendations are not sufficiently governed.
\item
  \texttt{GAP-DERMA-PATH-01} - \texttt{predicted\_pathology} is
  source-observed but its governed obligation status and
  diagnostic-authority boundary remain not established.
\end{itemize}

\subsection{MR-01 cumulative D3
disposition}\label{mr-01-cumulative-d3-disposition}

\textbf{PASS.} FR-01 and FR-02 provide the independently assessable
operational obligations currently justified under DEC-01 and DEC-02.
Remaining source meanings are preserved as gaps or routed evidence
rather than invented requirements.

\section{4. MR-02 - Indirizzamento
specialistico}\label{4-mr-02---indirizzamento-specialistico}

\subsection{Intent}\label{intent-1}

Determinare un\textquotesingle indicazione di destinazione specialistica
per il caso dermatologico, al fine di supportarne il successivo
instradamento verso l\textquotesingle assistenza appropriata.

\subsection{Status}\label{status}

\textbf{STOP AT MR.}

The source supports that DermaTriage includes specialist-oriented
output/indication, but does not sufficiently govern:

\begin{itemize}
\tightlist
\item
  specialist-domain vocabulary;
\item
  selection rule;
\item
  required inputs;
\item
  fallback behavior;
\item
  actual assignment/booking responsibility;
\item
  complete relationship with P-scale priority.
\end{itemize}

No Decision or FR is created merely to fill the hierarchy.

\section{5. MR-03 - Gestione della validazione clinica degli
esiti}\label{5-mr-03---gestione-della-validazione-clinica-degli-esiti}

\subsection{Intent}\label{intent-2}

Supportare la registrazione e l\textquotesingle utilizzo della
validazione o correzione effettuata da un professionista sanitario sugli
esiti prodotti da DermaTriage, mantenendo distinguibile il risultato
della revisione dall\textquotesingle esito originario.

\subsection{Responsibility boundary}\label{responsibility-boundary}

The healthcare professional owns the clinical judgment. DermaTriage owns
management, recording and correlation of the review result.

\subsection{DEC-03 - Separazione tra esito originario e revisione
clinica}\label{dec-03---separazione-tra-esito-originario-e-revisione-clinica}

\textbf{Decision kind:} SELECTABLE\\
\textbf{Status:} PASS

DermaTriage keeps the original system outcome distinguishable from the
subsequent clinical review result.

\subsubsection{FR-03 - Registrazione correlata della revisione
clinica}\label{fr-03---registrazione-correlata-della-revisione-clinica}

\textbf{Parent:} DEC-03\\
\textbf{Status:} PASS

\textbf{NC-03.1}\\
When a clinical validation/correction is supplied for a DermaTriage
outcome, DermaTriage MUST record the review result associated with the
original outcome.

\textbf{NC-03.2}\\
The recorded review result MUST distinguish confirmation from
correction.

\subsubsection{Preserved gaps}\label{preserved-gaps-1}

\begin{itemize}
\tightlist
\item
  \texttt{GAP-DERMA-REVIEW-CONTENT-01} - the complete correction/review
  content model is not specified.
\item
  \texttt{GAP-DERMA-REVIEW-LIFECYCLE-01} -
  overwrite/history/retention/finality semantics are not specified.
\item
  \texttt{GAP-DERMA-REVIEW-DISAGREEMENT-BINDING-01} - the normative
  relationship among correction, disagreement and the source-observed
  \texttt{agrees} state is not fully specified.
\end{itemize}

\subsection{MR-03 cumulative D3
disposition}\label{mr-03-cumulative-d3-disposition}

\textbf{PASS.}

\section{6. MR-04 - Adattamento controllato sulla base della revisione
clinica}\label{6-mr-04---adattamento-controllato-sulla-base-della-revisione-clinica}

\subsection{Intent}\label{intent-3}

Consentire l\textquotesingle evoluzione controllata del comportamento di
DermaTriage utilizzando evidenza derivata dalle validazioni e correzioni
cliniche accumulate, evitando che modifiche degradanti vengano accettate
senza controllo.

\subsection{Dependency}\label{dependency}

\texttt{MR-04\ dependsOn\ MR-03\ =\ PASS} because adaptation consumes
clinical-review meaning owned by MR-03.

\subsection{Decision family}\label{decision-family}

\begin{verbatim}
DEC-04  Accumulation-gated adaptation
DEC-05  Separate adaptation paths
DEC-06  Comparative acceptance
DEC-07  Asymmetric quality prioritization
DEC-08  Classification adaptation reversibility
DEC-09  Bounded recent evidence - prompt path
DEC-10  Disagreement as qualifying classifier evidence
DEC-11  Corrected P-scale as supervision source
\end{verbatim}

\subsection{DEC-04 - Attivazione dell\textquotesingle adattamento su
accumulo di evidence
clinica}\label{dec-04---attivazione-delladattamento-su-accumulo-di-evidence-clinica}

DermaTriage uses a non-immediate adaptation strategy: each adaptation
process is activated only when relevant accumulated clinical evidence
satisfies a governed accumulation condition for that process.

\subsubsection{FR-04 - Attivazione del percorso di evoluzione del
prompt}\label{fr-04---attivazione-del-percorso-di-evoluzione-del-prompt}

\textbf{Parent:} DEC-04\\
\textbf{Status:} PASS

\begin{quote}
Quando l\textquotesingle evidence di correzione clinica pertinente
accumulata raggiunge la soglia prevista per il percorso di evoluzione
del prompt, DermaTriage MUST attivare un ciclo di evoluzione del prompt.
\end{quote}

Semantic parameter: \texttt{PromptEvolutionThreshold=N}\\
Current documented binding: \texttt{N=10}\\
Exact fixed normative status: \texttt{NOT\ SPECIFIED}\\
Change authority: \texttt{NOT\ SPECIFIED}

\subsubsection{FR-05 - Attivazione del percorso di adattamento
classificatorio su accumulo di evidence di
disaccordo}\label{fr-05---attivazione-del-percorso-di-adattamento-classificatorio-su-accumulo-di-evidence-di-disaccordo}

\textbf{Parent:} DEC-04\\
\textbf{Status:} PASS

\begin{quote}
Quando l\textquotesingle evidence qualificante accumulata raggiunge la
soglia prevista per il percorso di adattamento classificatorio,
DermaTriage MUST attivare il relativo ciclo di adattamento.
\end{quote}

Semantic parameter: \texttt{ClassifierAdaptationThreshold=N}\\
Current documented binding: \texttt{N=50}\\
Exact fixed normative status: \texttt{NOT\ SPECIFIED}\\
Change authority: \texttt{NOT\ SPECIFIED}

\subsubsection{Preserved gaps under
DEC-04}\label{preserved-gaps-under-dec-04}

\begin{itemize}
\tightlist
\item
  \texttt{GAP-DERMA-ADAPT-COUNTING-01} - reset, deduplication, reuse,
  persistence, multiple-review and concurrency counting semantics are
  not specified.
\item
  \texttt{GAP-DERMA-ADAPT-THRESHOLD-AUTH-01} - governance/change
  authority of concrete threshold values is not specified.
\end{itemize}

\subsection{DEC-05 - Separazione dei percorsi di adattamento per
capability}\label{dec-05---separazione-dei-percorsi-di-adattamento-per-capability}

DermaTriage governs adaptation through distinct lifecycle paths for
different behavioral capabilities. Conditions or outcomes belonging to
one path do not automatically satisfy those of another path.

\subsubsection{FR-11 - Indipendenza dei lifecycle di
adattamento}\label{fr-11---indipendenza-dei-lifecycle-di-adattamento}

\textbf{Parent:} DEC-05\\
\textbf{Status:} PASS / STOP

\textbf{NC-11.1}\\
The satisfaction of an activation condition for one adaptation path MUST
NOT, by itself, be treated as satisfaction of the activation condition
of another path.

\textbf{NC-11.2}\\
Evidence qualifying for one adaptation path MUST NOT, by itself, be
treated as qualifying for another path.

\textbf{NC-11.3}\\
A qualification/progression result reached by one adaptation path MUST
NOT, by itself, be treated as an equivalent lifecycle result for another
path.

Shared data may be reused if it independently satisfies the governed
rules of each path; shared data does not imply shared lifecycle
semantics.

\subsection{DEC-09 - Uso di evidence recente e limitata nel percorso di
evoluzione del
prompt}\label{dec-09---uso-di-evidence-recente-e-limitata-nel-percorso-di-evoluzione-del-prompt}

DermaTriage uses bounded recent clinical-correction evidence rather than
unrestricted history for prompt evolution.

\subsubsection{FR-06 - Selezione dell\textquotesingle evidence recente
per l\textquotesingle evoluzione del
prompt}\label{fr-06---selezione-dellevidence-recente-per-levoluzione-del-prompt}

\textbf{Parent:} DEC-09\\
\textbf{Status:} PASS / STOP

\begin{quote}
Per ciascun ciclo di evoluzione del prompt, DermaTriage MUST costruire
l\textquotesingle evidence set utilizzando una finestra scorrevole e
limitata delle correzioni cliniche piu\textquotesingle{} recenti
pertinenti al percorso di evoluzione del prompt.
\end{quote}

Semantic parameter: \texttt{PromptEvidenceWindowSize=N}\\
Current documented binding: \texttt{N=20}\\
Exact fixed normative status: \texttt{NOT\ SPECIFIED}\\
Change authority: \texttt{NOT\ SPECIFIED}

Preserved gap: \texttt{GAP-DERMA-PROMPT-WINDOW-01} - ordering,
membership, underfill, deduplication, reuse/overlap and exact scope
semantics are not specified.

\subsection{DEC-10 - Qualificazione del disaccordo clinico come evidence
per l\textquotesingle adattamento
classificatorio}\label{dec-10---qualificazione-del-disaccordo-clinico-come-evidence-per-ladattamento-classificatorio}

For urgency-classification adaptation, DermaTriage treats clinical
reviews expressing disagreement with the relevant AI outcome as
qualifying evidence on the disagreement dimension.

\subsubsection{FR-07 - Qualificazione del disaccordo clinico come
evidence per l\textquotesingle adattamento
classificatorio}\label{fr-07---qualificazione-del-disaccordo-clinico-come-evidence-per-ladattamento-classificatorio}

\textbf{Parent:} DEC-10\\
\textbf{Status:} PASS / STOP

\begin{quote}
Nel determinare l\textquotesingle evidence clinica qualificante per il
percorso di adattamento della classificazione di urgenza, DermaTriage
MUST considerare qualificanti soltanto le revisioni cliniche che
esprimono disaccordo rispetto all\textquotesingle esito AI pertinente.
\end{quote}

Governed concept: \texttt{ClinicianDisagreement}\\
Source-observed current encoding: \texttt{agrees\ ==\ False}\\
The encoding is realization/data-state evidence, not the normative
semantic concept.

Confirmed consumption: \texttt{FR-07\ -\textgreater{}\ FR-05}.

\subsection{DEC-11 - Derivazione della supervision classificatoria dalla
priorita\textquotesingle{} clinicamente
corretta}\label{dec-11---derivazione-della-supervision-classificatoria-dalla-priorita-clinicamente-corretta}

For urgency-classification adaptation, DermaTriage uses the clinically
corrected P-scale priority as the governed source from which to derive
the classifier-supervision target.

\subsubsection{FR-08 - Derivazione del target di supervisione dalla
priorita\textquotesingle{} P-scale
corretta}\label{fr-08---derivazione-del-target-di-supervisione-dalla-priorita-p-scale-corretta}

\textbf{Parent:} DEC-11\\
\textbf{Status:} PASS / STOP

For classifier-adaptation cases with a clinically corrected P-scale
priority, DermaTriage MUST derive the supervision target according to:

\begin{verbatim}
P1 -> HIGH
P2 -> HIGH
P3 -> MEDIUM
P4 -> LOW
\end{verbatim}

The transformation is governed semantic meaning, not a replaceable
\texttt{N} parameter.

Preserved gap: \texttt{GAP-DERMA-SUPERVISION-INPUT-01} - missing,
invalid, out-of-domain or conflicting corrected-priority behavior is not
specified.

\subsection{DEC-06 - Accettazione comparativa
dell\textquotesingle adattamento della classificazione di
urgenza}\label{dec-06---accettazione-comparativa-delladattamento-della-classificazione-di-urgenza}

A candidate urgency-classification adaptation is not adopted merely
because it has been produced; it must first pass comparative evaluation
against the applicable reference under governed acceptance criteria.

\subsubsection{FR-09 - Qualificazione comparativa del candidato di
adattamento
classificatorio}\label{fr-09---qualificazione-comparativa-del-candidato-di-adattamento-classificatorio}

\textbf{Parent:} DEC-06\\
\textbf{Status:} PASS / STOP

\textbf{NC-09.1}\\
When a candidate urgency-classification adaptation is considered for
adoption, DermaTriage MUST evaluate it comparatively against the
applicable reference version using governed acceptance criteria.

\textbf{NC-09.2}\\
The candidate MUST be considered qualified for adoption only if all
applicable governed comparative acceptance criteria are satisfied.

\texttt{qualified\ for\ adoption} does not mean
\texttt{automatically\ deployed}.

Preserved gap: \texttt{GAP-DERMA-DEPLOY-01} - final promotion authority
and automaticity after qualification are not specified.

\subsection{DEC-07 - Prioritizzazione asimmetrica delle metriche di
qualita\textquotesingle{}}\label{dec-07---prioritizzazione-asimmetrica-delle-metriche-di-qualita}

DermaTriage governs an asymmetric acceptance policy in which sensitivity
and false-low performance may not worsen relative to the reference,
while overall accuracy may degrade only within a bounded tolerance.

\subsubsection{A5 disposition}\label{a5-disposition}

\textbf{NO NEW FR / ROUTE TO A7.}

This Decision governs additional quality properties of FR-09 rather than
a new autonomous operational capability.

Current specialization candidates to preserve for A7:

\begin{verbatim}
SensitivityNonDegradation
FalseLowNonDegradation
LimitedAccuracyDegradation
\end{verbatim}

Semantic parameter for the third property:
\texttt{AccuracyDegradationTolerance=T}\\
Current textual binding: \texttt{T=5\%}\\
Absolute-versus-relative interpretation: \texttt{NOT\ SPECIFIED}\\
Exact reference/evaluation population/change authority:
\texttt{NOT\ SPECIFIED}

Concrete non-security subtype design is outside the current thesis
specialization scope unless later explicitly reopened.

Preserved gap: \texttt{GAP-DERMA-ACCEPT-BINDING-01}.

\subsection{DEC-08 - Reversibilita\textquotesingle{}
dell\textquotesingle adattamento della classificazione in caso di
degrado}\label{dec-08---reversibilita-delladattamento-della-classificazione-in-caso-di-degrado}

An already adopted urgency-classification adaptation that manifests
material degradation according to governed post-adoption criteria must
be reversible by restoring a previous acceptable version or state.

\subsubsection{FR-10 - Revoca di un adattamento classificatorio
degradante}\label{fr-10---revoca-di-un-adattamento-classificatorio-degradante}

\textbf{Parent:} DEC-08\\
\textbf{Status:} PASS / STOP

\begin{quote}
Per un adattamento della classificazione di urgenza
gia\textquotesingle{} adottato che manifesti degrado materiale secondo i
criteri post-adoption governati, DermaTriage MUST supportarne la revoca
mediante il ripristino di una versione o stato precedente accettabile.
\end{quote}

The requirement does not establish that rollback must be automatic.

Semantic parameter: \texttt{RollbackAccuracyDegradationThreshold=R}\\
Current textual binding: \texttt{R=5\%}\\
\texttt{R} is not assumed identical to pre-adoption \texttt{T} merely
because the current literals are equal.

Preserved gap: \texttt{GAP-DERMA-ROLLBACK-BINDING-01} - unit, reference,
population, observation window, timing, automaticity, authorization and
exact rollback target remain insufficiently governed.

\subsection{MR-04 cumulative D3
disposition}\label{mr-04-cumulative-d3-disposition}

\textbf{PASS / CLOSED at A5.}

No MR-04 Decision remains without an A5 disposition:

\begin{longtable}[]{@{}ll@{}}
\toprule\noalign{}
Decision & A5 result \\
\midrule\noalign{}
\endhead
\bottomrule\noalign{}
\endlastfoot
DEC-04 & FR-04 + FR-05 PASS \\
DEC-05 & FR-11 PASS / STOP \\
DEC-06 & FR-09 PASS / STOP \\
DEC-07 & no FR; route to A7 \\
DEC-08 & FR-10 PASS / STOP \\
DEC-09 & FR-06 PASS / STOP \\
DEC-10 & FR-07 PASS / STOP \\
DEC-11 & FR-08 PASS / STOP \\
\end{longtable}

\section{7. Cross-requirement and cross-MR
relations}\label{7-cross-requirement-and-cross-mr-relations}

Confirmed:

\begin{verbatim}
MR-04 dependsOn MR-03
FR-07 -> FR-05   [qualifying disagreement evidence is accumulated by FR-05]
\end{verbatim}

Not asserted without additional evidence:

\begin{verbatim}
FR-01 -> FR-02
FR-07 -> FR-08
FR-08 -> FR-09
FR-09 -> FR-10
MR-04 dependsOn MR-01
\end{verbatim}

MR-04 may reference/consume the P-scale domain governed by MR-01 without
acquiring a second parent or automatically creating a macro dependency.

\section{8. Parameter and binding registry - current
candidate}\label{8-parameter-and-binding-registry---current-candidate}

\begin{longtable}[]{@{}llrl@{}}
\toprule\noalign{}
Semantic role & DDTA semantic form & Current source binding & Governance
status \\
\midrule\noalign{}
\endhead
\bottomrule\noalign{}
\endlastfoot
prompt adaptation accumulation threshold &
\texttt{PromptEvolutionThreshold=N} & 10 & exact value fixed status NOT
SPECIFIED \\
classifier adaptation accumulation threshold &
\texttt{ClassifierAdaptationThreshold=N} & 50 & exact value fixed status
NOT SPECIFIED \\
prompt recent-evidence window size & \texttt{PromptEvidenceWindowSize=N}
& 20 & exact value fixed status NOT SPECIFIED \\
pre-adoption accuracy degradation tolerance &
\texttt{AccuracyDegradationTolerance=T} & 5\% & quantitative binding
incomplete \\
post-adoption rollback accuracy degradation threshold &
\texttt{RollbackAccuracyDegradationThreshold=R} & 5\% & quantitative
binding incomplete \\
\end{longtable}

\texttt{T\ ==\ R} is not established.

\section{9. Information/data-state bridge ledger - current
candidate}\label{9-informationdata-state-bridge-ledger---current-candidate}

\begin{longtable}[]{@{}lll@{}}
\toprule\noalign{}
Governed semantic concept/reference & Source-observed/current binding &
Status \\
\midrule\noalign{}
\endhead
\bottomrule\noalign{}
\endlastfoot
\texttt{ClinicianDisagreement} & \texttt{agrees\ ==\ False} & data/state
encoding evidence; not normative field identity \\
\texttt{ClinicallyCorrectedOperationalTriagePriority} & source
review/correction data carrying P1-P4 & exact interface field binding
pending A10 \\
\texttt{ClassifierUrgencyTarget} & HIGH/MEDIUM/LOW; possible
enum/string/index realizations & semantic domain governed; concrete
encoding pending A10 \\
\texttt{PromptEvolutionThreshold} & current 10 & parameter binding \\
\texttt{ClassifierAdaptationThreshold} & current 50 & parameter
binding \\
\texttt{PromptEvidenceWindowSize} & current 20 & parameter binding \\
\end{longtable}

\section{10. Preserved project gap
register}\label{10-preserved-project-gap-register}

\begin{longtable}[]{@{}lll@{}}
\toprule\noalign{}
Gap & Area & Preserved meaning \\
\midrule\noalign{}
\endhead
\bottomrule\noalign{}
\endlastfoot
Diagnostic authority boundary & framing / MR-01 & definitive clinical
diagnosis authority not established \\
\texttt{GAP-DERMA-NOIMAGE-INPUT-01} & FR-01 & complete symptom input
contract not specified \\
\texttt{GAP-DERMA-PMAP-INPUT-01} & FR-02 & complete urgency/confidence
input semantics not specified \\
\texttt{GAP-DERMA-NOIMAGE-PMAP-BINDING-01} & MR-01 & no-image urgency
-\textgreater{} P-scale input binding not established \\
\texttt{GAP-DERMA-SLA-01} & MR-01 downstream & SLA trigger/outcome/owner
semantics not established \\
\texttt{GAP-DERMA-PATH-01} & MR-01 downstream & pathology output
obligation/authority not established \\
specialist selection semantics & MR-02 &
domain/rules/input/fallback/booking not specified \\
\texttt{GAP-DERMA-REVIEW-CONTENT-01} & FR-03 & complete review payload
not specified \\
\texttt{GAP-DERMA-REVIEW-LIFECYCLE-01} & FR-03 &
overwrite/history/retention/finality not specified \\
\texttt{GAP-DERMA-REVIEW-DISAGREEMENT-BINDING-01} & MR-03/MR-04 &
correction/disagreement/\texttt{agrees} relation incomplete \\
\texttt{GAP-DERMA-ADAPT-COUNTING-01} & FR-04/05 &
counting/reset/dedup/reuse semantics not specified \\
\texttt{GAP-DERMA-ADAPT-THRESHOLD-AUTH-01} & FR-04/05 & threshold change
authority not specified \\
\texttt{GAP-DERMA-PROMPT-WINDOW-01} & FR-06 & exact window
membership/lifecycle semantics incomplete \\
\texttt{GAP-DERMA-SUPERVISION-INPUT-01} & FR-08 & invalid/missing
corrected-priority behavior not specified \\
\texttt{GAP-DERMA-DEPLOY-01} & FR-09 lifecycle & final deployment
authority/automaticity not specified \\
\texttt{GAP-DERMA-ACCEPT-BINDING-01} & DEC-06/07 & complete quantitative
acceptance binding not specified \\
\texttt{GAP-DERMA-ROLLBACK-BINDING-01} & DEC-08 & complete rollback
trigger binding not specified \\
\end{longtable}

\section{11. Source-observed evidence pending A7/A8/A10
review}\label{11-source-observed-evidence-pending-a7a8a10-review}

The following source facts are preserved but \textbf{not yet promoted}
to governed Specialized/Security requirements or stable realization
bindings by this candidate:

\subsection{11.1 Non-security quality specialization evidence - pending
A7}\label{111-non-security-quality-specialization-evidence---pending-a7}

\begin{itemize}
\tightlist
\item
  sensitivity non-degradation relative to the applicable reference;
\item
  false-low non-degradation relative to the applicable reference;
\item
  bounded accuracy degradation;
\item
  source documents report a current \texttt{5\%} literal for accuracy
  tolerance.
\end{itemize}

\subsection{11.2 Security/privacy evidence - pending
A8}\label{112-securityprivacy-evidence---pending-a8}

Source documents currently describe:

\begin{itemize}
\tightlist
\item
  \texttt{X-API-Key} on protected DermaTriage
  administrative/B4-integrated operations;
\item
  B4 JWT authentication;
\item
  anonymized research patient identifiers using a \texttt{PAT\_XXXX}
  style;
\item
  no raw patient images stored in the RAG CSV;
\item
  B4 patient uploads processed in memory and not stored permanently in
  the research dataset.
\end{itemize}

These observations are \textbf{A8 candidate evidence}, not yet final
SecurityRequirements. Mechanisms such as X-API-Key/JWT are not
automatically the security property.

\subsection{11.3 Realization/configuration evidence - pending
A10}\label{113-realizationconfiguration-evidence---pending-a10}

Examples include:

\begin{itemize}
\tightlist
\item
  current named models and four-stage pipeline;
\item
  endpoint paths;
\item
  environment variables;
\item
  prompt storage path;
\item
  database filenames;
\item
  fine-tuning layers, learning rates, epochs and class weights;
\item
  model version/file identifiers;
\item
  current threshold tuning values not already promoted to governed
  semantics.
\end{itemize}

\section{12. Current gate state}\label{12-current-gate-state}

\begin{verbatim}
A0  Authority/source gate                  PASS
A1  Project framing                        PASS
A2  MacroRequirement / D1                  CUMULATIVE PASS
A3  Semantic sufficiency                   PASS for decomposed scope
A4  Decision / D2                          CUMULATIVE PASS / CLOSED
A5  FunctionalRequirement / D3             CUMULATIVE PASS / CLOSED
A6  Requirement split                      NEXT
A7  SpecializedRequirement                 NOT STARTED
A8  SecurityRequirement                    NOT STARTED
A9  Cross-MR downstream review             NOT STARTED
A10 Terminology/information/bindings        NOT STARTED
A11 Propagation/regression                  NOT STARTED
A12 Completeness/promotion readiness        NOT STARTED
BA   Accepted Base Analysis                 BLOCKED
\end{verbatim}

This candidate must not be promoted as the final DermaTriage governed
baseline before A6-A12 complete.

\end{document}
